\section*{第 \ref{ch:b3:planetary-interiors-surfaces} 章——行星内部与表面}

\begin{solution}{exo:b3:planetary-interiors-surfaces:1}
$\tfrac52C\bar\rho = 2.5\times0.35\times5429 = 4750$; $x^2 = (4750 -
3300)/(5429 - 3300) = 0.681$，$x = 0.825$，核心半径为 $0.825\times2440 =
\qty{2010}{km}$; $\rho_c = 3300 + 2129/0.825^3 = \qty{7090}{kg.m^{-3}}$.
\end{solution}

\begin{solution}{exo:b3:planetary-interiors-surfaces:2}
$m = 3500\times\tfrac43\pi(25)^3 = \qty{2.29e8}{kg}$; $E = \tfrac12 m v^2 =
\tfrac12\times\num{2.29e8}\times(\num{2.0e4})^2 = \qty{4.6e16}{J}$，或
$\num{4.6e16}/\num{4.2e15} = \qty{11}{Mt}$。由经验规则得出的撞击坑直径为
$20\times\qty{50}{m} = \qty{1}{km}$.
\end{solution}

\begin{solution}{exo:b3:planetary-interiors-surfaces:3}
$N_1 = 120/\num{100000} = \qty{1.2e-3}{km^{-2}}$，相对误差 $1/\sqrt{120}
= \qty{9}{\percent}$。在线性段上，$T = N_1/c = \num{1.2e-3}/\num{4.875e-4}
= \qty{2.46}{Ga}$（指数项 $\num{8e-7}$ 可忽略），且
$\sigma_T = T/\sqrt{120} = \qty{0.22}{Ga}$.
\end{solution}

\begin{solution}{exo:b3:planetary-interiors-surfaces:4}
面积为 $4\pi(\num{1.8215e6})^2 = \qty{4.17e13}{m^2}$；功率为 $2.5\times
\num{4.17e13} = \qty{1.04e14}{W} = \qty{104}{TW}$，是地球总功率的 2.2 倍，
来自表面积仅为地球十二分之一的表面，且 $2.5/0.092 = 27$ 倍于地球的
热流。
\end{solution}

\begin{solution}{exo:b3:planetary-interiors-surfaces:5}
(a)~$\bar\rho = \num{6.417e23}/(\tfrac43\pi(\num{3.3895e6})^3) =
\qty{3934}{kg.m^{-3}}$; $g = GM/R^2 = \qty{3.73}{m.s^{-2}}$; $\rho gR =
\qty{4.97e10}{Pa}$. (b)~$19\mu/(2\rho gR) = 19\times\num{6.7e10}/\num{9.94e10}
= 12.8$，$k_2 = 1.5/13.8 = 0.109$。(c)~实测值比
大 2.6 个标准差：火星的表现超过了地幔刚性相同的固体行星，
正如其内部含有液态核心时那样。
\end{solution}

\begin{solution}{exo:b3:planetary-interiors-surfaces:6}
(a)~月球：结晶岩中为 $4\times9.8/1.62 = \qty{24}{km}$，
沉积物中为 \qty{12}{km}；火星和水星：$4\times9.8/3.7 =
\qty{10.6}{km}$ 和 \qty{5.3}{km}。月球的 \qty{20}{km} 位于
两个预测值之间：该规律的精度在百分之几十以内。(b)~$Y =
\rho gD/3 = 2600\times9.8\times4000/3 = \qty{34}{MPa}$，这是破碎岩石而非完整晶体的
强度。
\end{solution}

\begin{solution}{exo:b3:planetary-interiors-surfaces:7}
火星：$\num{3.22e-14}(e^{6.93T} - 1) + \num{4.875e-4}\,T = \num{2.0e-3}$；
线性估计 $4.1$~Ga 远远偏大，
此处指数项很大；二分法给出 $T = \qty{3.34}{Ga}$（指数项
$\num{3.6e-4}$，线性项 $\num{1.63e-3}$）。月球函数给出
\qty{2.39}{Ga}。每个函数只针对一个天体标定：撞击
速率，以及同一撞击体所形成撞击坑的大小，在月球与
火星之间各不相同。
\end{solution}

\begin{solution}{exo:b3:planetary-interiors-surfaces:8}
(a)~$3.551181/1.769138 = 2.0073$ 和 $7.154553/3.551181 = 2.0147$，接近
$2$。(b)~$0.565247 - 0.844789 + 0.279542 = 0.000000$，每天，精确到
$\num{1e-7}$：三者的平均运动被锁定，形成一种共振，
它保持了三颗卫星在每次合冲时的相对位置。(c)~拉力在
轨道上的相同位置重复出现，因此它们叠加起来，而不是相互抵消；
它们补充偏心率的速度与潮汐使其衰减的速度相当。
\end{solution}

\begin{solution}{exo:b3:planetary-interiors-surfaces:9}
$n = 2\pi/(1.769138\times86400) = \qty{4.111e-5}{s^{-1}}$, $nR =
\qty{74.9}{m.s^{-1}}$, $(nR)^5/G = \qty{3.53e19}{W}$;
$\tfrac{21}{2}e^2(nR)^5/G = 10.5\times\num{1.6e-5}\times\num{3.53e19} =
\qty{5.9e15}{W}$（每单位 $k_2/Q$）。则 $k_2/Q =
\num{1.04e14}/\num{5.9e15} = 0.018$.
\end{solution}

\begin{solution}{exo:b3:planetary-interiors-surfaces:10}
(a)~$x^2 = (\tfrac52C\bar\rho - \rho_m)/(\bar\rho - \rho_m)$ 关于
$C$ 呈线性，其斜率为 $\tfrac52\bar\rho/(\bar\rho - \rho_m)$。(b)~该斜率为
$2.5\times3340/52 = 161$：$x^2$ 改变 $0.161$，对应 $0.001$ 的 $C$ 变化。
当 $C = 0.395$ 时，$x^2 = 0.036 + 0.161 = 0.197$，$x = 0.44$，$\rho_c =
\qty{3880}{kg.m^{-3}}$；当 $C = 0.393$ 时，$x^2 < 0$：根本不存在核心。
(c)~在 $C = \tfrac25$ 附近，转动惯量对小核几乎不敏感：
其第三位有效数字决定了存在一个大而轻的核心，还是根本没有核心。
还需要借助地震波、潮汐和磁场。
\end{solution}

\begin{solution}{exo:b3:planetary-interiors-surfaces:11}
(a)~$\sigma_T = N_1/(10\,N_1')$：在 \qty{1.0}{Ga} 时，$N_1 \approx N_1' \approx
\num{4.875e-4}$，$\sigma_T = \qty{0.10}{Ga}$；在 \qty{3.8}{Ga} 时，$N_1 =
\num{1.07e-2}$ 和 $N_1' = \num{6.15e-2}$，$\sigma_T = \qty{0.017}{Ga}$，
小六倍。(b)~若 $N_1 = cT$，$N_1' = c$ 且 $\sigma_T =
cT/(\sqrt n\,c) = T/\sqrt n$。(c)~年代函数本身是
不确定的：它依赖于月球上少数已定年的地点，以及由月球向火星的
标度关系，这些误差无论计数多少都无法减小。
\end{solution}

\begin{solution}{exo:b3:planetary-interiors-surfaces:12}
(a)~比值为 $(n_ER_E/n_IR_I)^5(e_E/e_I)^2 = (0.4982\times0.8569)^5
\times(0.009/0.004)^2 = 0.0142\times5.06 = 0.072$：约
\qty{7.5e12}{W}，即 $\num{7.5e12}/(4\pi(\num{1.5608e6})^2) =
\qty{0.24}{W.m^{-2}}$。(b)~冰下海洋使冰壳与
内部脱耦：冰壳的挠曲远大于固体卫星，其冰
耗散方式也与木卫一炽热的岩石不同，因此 $k_2$ 和 $Q$ 都会改变。
\end{solution}

\begin{solution}{pb:b3:planetary-interiors-surfaces:1}
\textbf{1.} $\bar\rho = \num{6.4169e23}/(\tfrac43\pi(\num{3.3895e6})^3) =
\qty{3934}{kg.m^{-3}}$.

\textbf{2.} $\tfrac52C\bar\rho = 3600$; $x^2 = (3600 - 3500)/(3934 - 3500) =
0.230$，$x = 0.479$，核心半径为 \qty{1620}{km}；$\rho_c = 3500 +
434/0.479^3 = \qty{7450}{kg.m^{-3}}$.

\textbf{3.} $x = 1830/3389.5 = 0.540$；$\rho_m = (3600 - 0.2915\times
3934)/(1 - 0.2915) = \qty{3460}{kg.m^{-3}}$，$\rho_c = 3462 +
472/0.1574 = \qty{6460}{kg.m^{-3}}$。

\textbf{4.} 比地球低约三分之一。火星核心的压强远低于地球，
因此其中的金属压缩程度较小；核心如此之轻，
说明铁中混有密度小得多的元素。

\textbf{5.} \omterm{def:b3:planetary-interiors-surfaces:love-number}{勒夫数}（2003 年）比 InSight（2021 年）的结果早得多：它已表明
至少外核是液态的，并将核心半径限定在 \qty{1520}{km} 至 \qty{1840}{km} 之间，
这一范围包含
了地震学的数值。

\textbf{6.} 体积占比为 $1 - (1 - 20/3389.5)^3 = \qty{1.8}{\percent}$ 至 $1 - (1 -
37/3389.5)^3 = \qty{3.2}{\percent}$，质量占比更小：
该层比地幔轻，故使有效地幔密度略有降低，
其效应小于模型的不确定度。

\textbf{7.} 从 \qty{1}{km} 起：54；从 1.41 起：16；从 2 起：10；从 2.83 起：5；
从 4、5.66、8 和 11.3 起：2；从 16 起：1。

\textbf{8.} $N = 54/80000 = \qty{6.75e-4}{km^{-2}}$，$\sigma = \sqrt{54}/80000
= \qty{0.92e-4}{km^{-2}}$（\qty{14}{\percent}）。

\textbf{9.} 预期数：从 \qty{1.41}{km} 起为 $54\times0.355 = 19$ 个撞击坑，
从 \qty{2}{km} 起为 $54\times0.147 = 7.9$ 个；实际计数为 16 和 10，
均在一个标准差之内（$\sqrt{19} = 4.4$，$\sqrt{7.9} = 2.8$）：一致。

\textbf{10.} 不是：该值低于 $0.072/\num{6.75e-4} = 107$ 倍。
撞击坑计数所记录的是时间。

\textbf{11.} 在 $\log D$ 中取等宽的区间，
可使各区间均匀覆盖整个范围，并在陡峭的分布中保有数量相近的撞击坑。

\textbf{12.} 它们是次生撞击坑，由从大型原生撞击坑中抛出的岩块形成：
它们并非来自太空的独立撞击，
若计入，将增加年代函数所不计的撞击坑。

\textbf{13.} $T = \num{6.75e-4}/\num{4.875e-4} = \qty{1.385}{Ga}$；指数项
$\num{3.22e-14}(e^{9.6} - 1) = \num{4.7e-10}$ 仅为
$N_1$ 的百万分之一。

\textbf{14.} $\sigma_T = T/\sqrt{54} = \qty{0.19}{Ga}$：$1.38 \pm 0.19$~Ga。

\textbf{15.} $\num{6.75e-4}/\num{8.38e-4} = \qty{0.81}{Ga}$，年轻了
将近一半。

\textbf{16.} $31 + 12 + 7 + 2 + 0 + 2 + 2 = 56$ 个撞击坑；$N(\qty{2}{km}) =
56/120000 = \qty{4.67e-4}{km^{-2}}$；$N_1 = \num{4.67e-4}/0.1473 =
\qty{3.17e-3}{km^{-2}}$。对 $\num{3.22e-14}(e^{6.93T} - 1) +
\num{4.875e-4}\,T = \num{3.17e-3}$ 用二分法求解，得 $T = \qty{3.54}{Ga}$（指数项
\num{1.45e-3}，线性项 \num{1.73e-3}）。

\textbf{17.} $N_1' = \num{3.22e-14}\times6.93\,e^{6.93\times3.54} +
\num{4.875e-4} = \num{1.05e-2}$；$\sigma_T = \num{3.17e-3}/(\sqrt{56}\times
\num{1.05e-2}) = \qty{0.040}{Ga}$。

\textbf{18.} 平原 A：\qty{14}{\percent}；平原 B：\qty{1.1}{\percent}。在
年代学曲线陡峭的早期部分，撞击坑密度的较大变化
只对应年龄的较小变化；在线性部分，年龄
的不确定度与计数本身的不确定度相当。

\textbf{19.} $120000\times\num{3.17e-3}\times(1 - 0.1473) = 324$；它携带
$48 + 8 = 56$ 个。

\textbf{20.} $N_1 = 56/(120000\times0.8527) = \qty{5.47e-4}{km^{-2}}$，
$T_r = \num{5.47e-4}/\num{4.875e-4} = \qty{1.12}{Ga}$。

\textbf{21.} $\sigma = 1.12/\sqrt{56} = \qty{0.15}{Ga}$。

\textbf{22.} 一个 \qty{2}{km} 的撞击坑深 \qty{400}{m}：熔岩
至少有这么厚。较大的撞击坑更深，坑缘也更高；
它们曾被部分填埋，至今仍清晰可见。

\textbf{23.} $1.38 \pm 0.19$ 与 $1.12 \pm 0.15$~Ga 相差 \qty{0.26}{Ga}，
约为一个合成标准差（$\sqrt{0.19^2 + 0.15^2} = 0.24$）：
计数结果允许存在单一事件，但不能证明这一点。

\textbf{24.} 年轻的平原 A 的模式年龄为
\textbf{14~亿年}（$1.38 \pm 0.19$~Ga）。
\end{solution}
